\subsection{The rank-four good residual lattice}
\label{subsec:sr-four-good-residual-lattice}

We supply the lattice producer preceding the Case II residual
chain. It is over the original rational group, before passing
to the auxiliary real tower.
Every value of the principal \(\psi\) has \(2\)-power order.
Thus the source set
\(S_b=\{p\ne\ell:\ell\mid\#\psi(I_p)\}\) is empty.

\paragraph{Normalization and the two permitted labels.}
Start with an actual integral Case II representation and its
ordinary symplectic flag, over a complete DVR model with finite
residue field. The representation reduces to the prescribed
diagonal characters. At the ordinary local group the characters
with labels \(0,2\) are separated from those with labels \(1,3\),
by the actual separator in the data. The first flag line
therefore reduces into the \(0,2\) coordinate plane.

Its primitive vector has either a unit \(0\) coordinate or a
unit \(2\) coordinate. In the second case interchange \(0,2\)
and simultaneously \(1,3\), using the symplectic permutation.
A unit \(2\) coordinate says that the first and third
characters agree on the whole ordinary local group: apply
each local group element to this flag vector and compare
that coordinate. Thus the interchange is permitted exactly
when \(\bar\Psi\) is trivial there, and preserves the
required local order. These are the two actual tower labels.

The first flag minor is now a unit, and so is the second.
Split the two separated residual sectors. Modulo the first
line the second flag vector lies in the other sector.
Orthogonality gives \(a_0b_3-a_2b_1=0\).
Since \(a_0\) is a unit, primitivity makes \(b_1\) a unit;
the second minor is \(a_0b_1\).
Subtracting a multiple of the first vector to obtain this
sector decomposition does not change the minor.

Choose the tower's \(g_0\), with four distinct residual
eigenvalues. Hensel spectral projectors diagonalize it by
a matrix congruent to \(1\). Its symplectic form pairs only
opposite eigenlines: their eigenvalue products are the
multiplier and residual separation excludes the other pairs.
Rescaling the paired eigenvectors gives a symplectic basis.
The flag minors stay units; lower/upper Gaussian factorization
gives its unique lower unipotent representative \(n\).
Its reduction has only the \(-\psi_{\rm root}\) coordinate,
since the other lower coordinates join the separated sectors.
Consequently
\[
 r(g_0)\text{ diagonal},\quad
 r(g)_{ij}\in\mathfrak m\ (i\ne j),\quad
                 n_{10},n_{21},n_{30}\in\mathfrak m,       \tag{R4V1}
\]
while \(n_{20}\) may be a unit.

An irreducible representation has every root among its entries:
a matrix entry missing on all group elements is missing in
their span, whereas irreducibility over the algebraic closure
makes that span the full matrix algebra.
For a nonsplit two-block point with irreducible outer rank-two
block, both \(\psi_{\rm root}\) and \(-\psi_{\rm root}\)
occur in that block. At least one of
\(\beta_{\rm root},-\alpha_{\rm root},-\theta_{\rm root}\)
occurs across the blocks, since otherwise it is block diagonal
and split. These statements survive normalization:
the permutation preserves the two blocks and the spectral
projectors may be chosen within the block parabolic.
Thus in either source \(\mathrm{Star}\) case we have
\[
 \operatorname{Ent}(\psi_{\rm root}),\
 \operatorname{Ent}(-\psi_{\rm root}),\
 \operatorname{Ent}(\beta_{\rm root})\ \text{or}\
 \operatorname{Ent}(-\alpha_{\rm root})\ \text{or}\
 \operatorname{Ent}(-\theta_{\rm root}).                   \tag{R4V2}
\]

\paragraph{The explicit valuation vertex.}
Our root coordinates are exactly
\[
 \alpha=(1,0),\quad\psi=(0,1),\quad
 \beta=(-1,1),\quad\theta=(1,1);
\]
the four matrix weights are
\(v_0=(1,1),v_1=(0,1),v_2=(1,0),v_3=(0,0)\),
and an \(ij\) entry has weight \(v_i-v_j\).
These are the actual \(\mathrm{lamW},\mathrm{muW}\)
coordinates in the frozen source, including its half-root
normalization, not the coordinates \(2e_1,2e_2\).
Normalize the original DVR valuation by \(v(\varpi)=1\).
For each root \(w\) include the valuations of all nonzero
matrix entries of weight \(w\); for each negative root also
include its independent flag coordinate
\(n_{10},n_{21},n_{20},n_{30}\).
Only their minima matter. Every finite minimum is attained,
and all these valuations are positive except possibly
the \(-\psi_{\rm root}\) flag valuation.

For integers \(A,B\) impose the finitely many inequalities
\[
                       w_1A+w_2B+2k\ge0                  \tag{R4V3}
\]
for these minima \(k\).
The feasible set contains \((0,0)\).
Its \(B\) coordinate is bounded above by (R4V2) at
\(-\psi_{\rm root}\).
For fixed \(B\), its \(A\) coordinate is bounded above by
one of the three cross-block roots in (R4V2), all with
first coordinate \(-1\).
Choose \(B\) maximal and then \(A\) maximal.
Let \(\mathcal S\) be the tight roots. Its possibilities are
\[
\begin{array}{c|c|c}
 &\text{necessarily tight}&\text{set containing }\mathcal S\\ \hline
1&\beta,-\psi&\{\beta,-\alpha,-\theta,-\psi\}\\
2&-\alpha,-\psi&\{-\alpha,-\psi,-\theta\}\\
3&-\theta,-\psi&\{-\theta,-\psi\}\\
4&-\alpha,-\beta&\{-\alpha,-\beta,-\psi,-\theta\}\\
5&-\theta,-\beta&\{-\theta,-\psi,-\beta\}\\
6&-\theta,\alpha&\{-\theta,-\psi,-\beta,\alpha\}.
\end{array}                                               \tag{R4V4}
\]
Root subscripts are suppressed in this table.

Here is the finite verification.
Maximality under \((1,0)\) and \((0,1)\) gives
\[
 T_{-\alpha}\vee T_\beta\vee T_{-\theta},\qquad
 T_{-\psi}\vee T_{-\beta}\vee T_{-\theta}.
\]
These tight equations imply \(A+B\) even: tight \(-\alpha\)
fixes \(A\) even, tight \(-\psi\) fixes \(B\) even,
and either long-root equation fixes \(A+B\) or \(A-B\) even.
The two disjunctions give the assertion in every combination.
The moves \((-1,1)\) and \((1,1)\) then give
\[
 T_{-\psi}\vee T_\alpha\vee T_{-\beta},\qquad
 T_{-\psi}\vee T_{-\alpha}\vee T_{-\theta}.
\]
A possible decrease by two cannot cross a positive odd slack,
since these long-root slacks are even.
For tight \(w\ne-\psi\), the original valuation is positive,
so \(w_1A+w_2B<0\); tight \(-\psi\) gives \(B\ge0\).
Tight \(\alpha,-\alpha,\beta,-\beta,\theta,-\theta,\psi\)
therefore imply, respectively,
\[
 A<0,\ A>0,\ B<A,\ A<B,\ A+B<0,\ A+B>0,\ B<0.
\]
Separate whether \(T_{-\psi}\) holds, then whether \(A<0\),
then whether \(T_{-\beta}\) holds. The four disjunctions and
these inequalities give exactly the six rows of (R4V4).
This checks all eight roots.

Adjoin \(p\) with \(p^2=\varpi\).
The grading action with parameters \(p^A,p^B\) is conjugation by
\(\operatorname{diag}(p^{A+B},p^B,p^A,1)\), a symplectic
similitude.
Since \(A+B\) is even, remove its scalar factor to use the
actual symplectic matrix
\[
 T=\operatorname{diag}\bigl(
 p^{(A+B)/2},p^{(B-A)/2},
 p^{(A-B)/2},p^{-(A+B)/2}\bigr),\qquad T^{\mathsf t}JT=J .
\]
Every displayed exponent is an integer. Conjugation
\(r\mapsto TrT^{-1}\), and likewise for \(n\), multiplies
an entry of weight \(w\) by \(p^{w_1A+w_2B}\).
An entry of weight \(w\) and valuation \(k\) has new valuation
\((w_1A+w_2B)/2+k\).
Thus every matrix entry and independent flag coordinate is
integral. The remaining lower flag entries are
\[
               n_{32}=-n_{10},\qquad
               n_{31}=n_{20}-n_{10}n_{21},
\]
so are integral too. The roots occurring in the reduction,
in either a matrix entry or independent flag coordinate,
are precisely \(\mathcal S\).
This constructs an actual valuation vertex over a finite
quadratic DVR extension.

\paragraph{The two implications eliminating the other vertices.}
For every odd root character \(\gamma\), a normalized residual
cocycle killed on the required inertia kernels and locally a
coboundary at \(\ell\) is zero.
Indeed Proposition~\ref{prop:sr-actual-corx} makes it a global
coboundary, since \(S_b=\varnothing\).
Its normalization at \(g_0\), where \(\gamma(g_0)\ne1\),
makes the global correcting scalar zero.
The local correcting scalar is zero too: the actual local
separator has \(\gamma(\sigma_\ell)\ne1\).
For extension-field coefficients the same assertion follows
by finite-field linear coordinate projections.

Consequently, if an odd-root entry \(r_{ij}\) has no other
contributing products in its multiplication equation, and
\[
 r(g)_{ij}=m_v(\bar\chi_j(g)-\bar\chi_i(g))
                                  \quad(g\in G_v),
\]
then \(r_{ij}=0\) and all \(m_v=0\):
the quotient \(r_{ij}/\bar\chi_j\) is precisely that normalized
root cocycle.

The ordinary first-column identity is
\[
 r(g)_{i0}+\sum_{j>0}r(g)_{ij}n_{j0}
                                =n_{i0}\bar\chi_0(g).     \tag{R4V5}
\]
If \(-\alpha\notin\mathcal S\), then
\(r_{10}=r_{32}=n_{10}=0\).
For \(i=3\), (R4V5) makes \(r_{30}\) the local coboundary
corrected by \(n_{30}\). Its other multiplication products
vanish. The preceding vanishing gives \(r_{30}=n_{30}=0\).
Hence
\[
 -\alpha\notin\mathcal S\Longrightarrow-\theta\notin\mathcal S.
                                                        \tag{R4V6}
\]
If \(\beta,\psi\notin\mathcal S\), then \(r_{12}=r_{13}=0\).
For \(i=1\), (R4V5) makes \(r_{10}\) the local coboundary
corrected by \(n_{10}\); its other products vanish.
Thus \(r_{10}=n_{10}=0\), and symplectic inversion gives
\(r_{32}=0\). Therefore
\[
 \beta,\psi\notin\mathcal S\Longrightarrow-\alpha\notin\mathcal S.
                                                        \tag{R4V7}
\]
Rows \(2,4\) contradict (R4V7), and rows \(3,5,6\)
contradict (R4V6). Only row \(1\) remains.

\paragraph{The actual good residual point and restriction.}
The reduced representation has support \(N_P\).
Its lower flag has support \(Q\), since
\(-\beta\notin\mathcal S\) gives \(n_{21}=0\).
The tight positive \(\beta\) root must be a matrix entry,
so \(r_{12}\ne0\).

If an entry at \(-\psi\) is nonzero, this is the first
residual alternative.
Otherwise tightness there gives \(n_{20}\ne0\), and the
global support is \(V\).
The \(i=2\) flag equation gives
\(\bar\chi_0|_{G_\ell}=\bar\chi_2|_{G_\ell}\),
which in either labeling is \(\bar\Psi|_{G_\ell}=1\).
The local \(12\) cocycle is not a coboundary:
otherwise its other products vanish in support \(N_P\),
and the odd-root vanishing would give \(r_{12}=0\).
This proves the full partially split alternative, including
its noncoboundary condition, hence the exact
\(\mathrm{LatShape}\).

Restriction to the real abelian tower gives the good point
required at the start of
Section~\ref{subsec:sr-four-case-II-residual-links}.
A normalized \(20\) or \(12\) cocycle cannot vanish on that
whole normal subgroup unless it vanishes globally:
conjugate a tower element at which its two characters differ
and use the cocycle equation.
The tower's regular \(g_0\) is such an element.
Both nonzero entries of the first alternative therefore
survive.
In the partially split alternative the transported \(20\)
flag coordinate remains nonzero by (R4L7).
If the restricted \(10\) entry vanished, that same
normal-subgroup argument would make it vanish globally.
Then the \(i=1\) flag equation, divided by \(n_{20}\ne0\),
would make the local \(12\) entry a coboundary, with correcting
scalar \(n_{10}/n_{20}\), a contradiction.
Thus the restricted point has support \(V\), every \(20\)
flag coordinate nonzero, and a nonzero \(10\) entry.
This is precisely the other good alternative used in the paths.

Conjugation preserves the multiplier, the ordered inertial
characters, the killed-inertia conditions and characteristic
polynomials.
Diagonalization, the square-root extension, torus powers and
flag elimination all take place over a finite extension of
the original complete DVR model, with finite residue field.
This proves the actual good residual lattice producer.

The frozen-source normalizations are
\srcid{FlagsAndDiagonals.Choice}, \srcid{ExNorm},
\srcid{StarConj}, and \srcid{CompactIntegralFields.Labels}.
The valuation calculation is
\srcid{RootSystemLinearProgramCases.h7} and
\srcid{AdicLattices.Vertex}.
Equations (R4V6)--(R4V7) expand
\srcid{OddVan}, \srcid{EntVan}, and \srcid{NotTight};
the final shape and restriction expand \srcid{LowCase},
\srcid{LatSix.lat}, \srcid{RestrQ}, and \srcid{FlagFacts}.
